\chapter{Les sources, côte à côte}\label{ch:tableau}

Les constats des chapitres précédents, résumés. « Cadre » signifie que le serveur laisse
afficher ses documents dans une page d'un autre site ; « CORS », qu'il laisse une page d'un
autre site lire ses fichiers dans le navigateur. Les constats sont datés dans leur
chapitre.

\begin{center}
\small\sffamily
\renewcommand{\arraystretch}{1.25}
\begin{tabular}{@{}p{2.3cm}p{2.0cm}p{2.5cm}p{3.6cm}@{}}
\textbf{Source} & \textbf{Format} & \textbf{Affichage} & \textbf{Droits et décision} \\ \hline
Montpellier & un PDF par dossier & cadre accepté depuis le 4 novembre 2026 (relais mis de côté) &
  protégé jusqu'en 2084 ; publié en démonstration, filigrané, retirable, sans DOI \\
Bourbaki (AHP) & un PDF par document, API Omeka & site en accès restreint (4 novembre 2026) &
  l'ACNB s'oppose à l'usage par IA ; travail suspendu, transcriptions privées \\
Quillen (Clay) & PDF & cadre accepté & protégé ; transcriptions privées, rien de plus
  transcrit \\
IHÉS & JPEG, TIFF, XML & vues de l'IHÉS dans le volet & correspondance protégée ; à
  confirmer avec l'IHÉS \\
BnF, Lasserre & pas en ligne & consultation en salle, sur déclaration & rien publié sans la BnF
  ni les héritiers \\
Buffalo & enregistrements, mis en ligne par l'Istituto Grothendieck & liens vers la vidéo &
  œuvre orale protégée ; un index, pas de transcription \\
Gallica & IIIF & CORS, pas de relais & domaine public ; conditions de Gallica ; CC0 sur ce
  qui est au projet \\
Leibniz (GWLB) & IIIF & CORS, pas de relais & domaine public ; \emph{Public Domain Mark}
  sur les manifestes relevés ; transcriptions d'autrui à respecter \\
\end{tabular}
\end{center}
